\subsection{An exact quadratic detector for the same-discriminant product}

The result in this subsection is a reformulation with an explicit integral
normalization.  It does not assert that the detector has a unit value,
and it does not supply the remaining auxiliary field.

\begin{proposition}[KT-EISENSTEIN-QUADRATIC-DETECTOR]
Let $p\geq31$ be prime with $p\equiv3\pmod4$, and put
\[
 b=(p+3)/2,\qquad k=b+3=(p+9)/2.
\]
Let $D<0$ be an odd fundamental discriminant, $d=|D|$, $p\nmid d$,
and write $\beta_j(\chi_D)=B_{j,\chi_D}/j$.  Normalize the level-one
Eisenstein series by
\[
 G_k(z)=-\frac{B_k}{2k}+\sum_{n\geq1}\sigma_{k-1}(n)e^{2\pi inz}.
\]
For the standard Gauss sum $\tau(\chi_D)$ one has the exact identity
\begin{equation}
 \frac{\tau(\chi_D)L(G_k\otimes\chi_D,3)}{(2\pi i)^3}
       =\frac{\beta_3(\chi_D)\beta_b(\chi_D)}{4d^2}.
 \tag{QD.1}
\end{equation}
Both Bernoulli quotients on the right are $p$-integral.

Choose a prime $q\equiv1\pmod p$, and set
\[
 E_{k,q}(z)=G_k(z)-q^{k-1}G_k(qz).
\]
This is a normalized Eisenstein eigenform of level $q$ whose constant
terms at every cusp vanish modulo $p$.  Consequently it has the
Heumann--Vatsal modular symbol $M_{k,q}$ modulo $p$.  For every $D$ as
above with $q\nmid d$ and $\chi_D(q)=-1$, the following coefficient is
\begin{equation}
 [X^2Y^{k-4}](M_{k,q}|R_{\chi_D})(\{\infty\}-\{0\})
       =-\frac{15}{8}\,\beta_3(\chi_D)\beta_b(\chi_D)
          \quad\hbox{in }\mathbb F_p.
 \tag{QD.2}
\end{equation}
In particular, this specified coefficient is nonzero if and only if both
required Bernoulli quotients are units.
\end{proposition}

\begin{proof}
For $\Re s>k$, the divisor-sum expansion gives
\[
 L(G_k\otimes\chi_D,s)=L(s,\chi_D)L(s-k+1,\chi_D),
\]
and continuation gives the identity at $s=3$.  The ordinary Bernoulli
Fourier series, followed by the primitive Gauss-sum identity, gives
\[
 \begin{split}
 B_{3,\chi_D}
 &=d^2\sum_{a=1}^{d}\chi_D(a)B_3(a/d)\\
 &=-\frac{6d^2}{(2\pi i)^3}
       \sum_{n\ne0}\frac{1}{n^3}
           \sum_{a=1}^{d}\chi_D(a)e^{2\pi ina/d}\\
 &=-\frac{12d^2\tau(\chi_D)}{(2\pi i)^3}L(3,\chi_D).
 \end{split}
\]
The Fourier series is absolutely convergent, and the last equality uses
the odd parity of $\chi_D$.  Since $L(1-b,\chi_D)=-\beta_b(\chi_D)$,
this proves (QD.1), including its sign and conductor power.
The Bernoulli-polynomial expression for $B_{j,\chi_D}$ shows its
$p$-integrality for $j=3,b$: here $j\leq p-2$, all ordinary Bernoulli
numbers of indices at most $j$ are $p$-integral, $p\nmid d$, and
$p\nmid j$.

Put $c=-B_k/(2k)$.  The inequalities $k<p-1$ and $p>k$ hold in the
stated range, so $c$ is $p$-integral.  The constant of $E_{k,q}$ at
infinity is $c(1-q^{k-1})$.  With the determinant-normalized slash
operator and $W_q=\left(\begin{smallmatrix}0&-1\\q&0\end{smallmatrix}\right)$,
\[
 G_k|_kW_q=q^{k/2}G_k(qz),\qquad
 G_k(qz)|_kW_q=q^{-k/2}G_k(z).
\]
Thus its constant at zero is $c q^{k/2-1}(q-1)$.  Both constants
vanish modulo $p$.  These are the two $\Gamma_0(q)$ cusp classes;
refining to $\Gamma_1(q)$ changes only the equivalent cusp
normalizations, with $p$-unit width factors.  The positive Fourier
coefficients are those of the normalized Eisenstein series attached to
the trivial character and the principal character modulo $q$.
In particular $a_1=1$, $U_qE_{k,q}=E_{k,q}$, and the good-prime
eigenvalues are $1+r^{k-1}$.

Theorem 3.23 of Heumann--Vatsal therefore supplies the integral modular
symbol modulo $p$.  Its construction is used here as a published input.
Twisting the displayed stabilization gives
\[
 \frac{\tau(\chi_D)L(E_{k,q}\otimes\chi_D,3)}{(2\pi i)^3}
  =(1-\chi_D(q)q^{k-4})
        \frac{\beta_3(\chi_D)\beta_b(\chi_D)}{4d^2}.
\]
Their Corollary 4.7 says that the coefficient in (QD.2), before
reduction, is the preceding value times
$-\binom{k-2}{2}2!d^2=-(k-2)(k-3)d^2$.
Now $q\equiv1$, $\chi_D(q)=-1$, $k-2\equiv5/2$, and
$k-3\equiv3/2$ modulo $p$.  The resulting scalar is $-15/8$.
It is a unit because $p\geq31$, proving the final equivalence.
\end{proof}

\paragraph{Exact source scope.}
J.~Heumann and V.~Vatsal, \emph{Modular symbols, Eisenstein series,
and congruences}, J. Th\'eor. Nombres Bordeaux \textbf{26} (2014),
709--756, DOI 10.5802/jtnb.886: the definition and hypotheses are on
pp.~734--735, Theorem 3.23 and its proof on pp.~736--737, and
Theorem 4.6 with its complete proof and Corollary 4.7 are on
pp.~742--745.  The primary PDF is
\url{https://www.numdam.org/item/JTNB_2014__26_3_709_0.pdf}.
Theorem 5.1 and its full proof on pp.~747--748 show that this modular
symbol is nonzero.  Their proof obtains a unit using characters with
growing auxiliary odd-prime-power order, and chooses the $j=0$
coefficient.  It does not imply (QD.2) is nonzero for any quadratic
character.  Formula (QD.1) was derived directly above; no odd-$j$
constant from their Proposition 6.1 is needed.

\paragraph{What a quadratic detector would still have to prove.}
The necessary nonzero evaluation is the particular coefficient (QD.2),
with an odd negative fundamental $D$, $p q\nmid D$, and
$\chi_D(q)=-1$.  A merely nonzero value in a different polynomial
component does not imply this conclusion.  The third auxiliary
condition $p\nmid h(D)$ must also be supplied.  It follows from the
already proved small-discriminant class bound if $d<4p$; it is not
automatic for an unbounded discriminant.  Thus neither full-symbol
primitivity nor unrestricted complex quadratic-twist nonvanishing is
the missing assertion.
